\documentclass[10pt,a4paper]{amsart}
\usepackage[margin=19mm]{geometry}
\usepackage{amsmath,amssymb,mathtools}
\usepackage[scr=rsfs,cal=euler]{mathalfa}
\usepackage{xfrac}
\usepackage[unicode,hidelinks,bookmarks=false]{hyperref}
\newcommand{\ie}{i.e.\ }
\newcommand{\cf}{cf.\ }
\newcommand{\resp}{respectively\ }
\newcommand{\Subsection}{\subsection}
\providecommand{\nobreakdash}{\nobreak\hbox{-}}
\setlength{\emergencystretch}{2em}
\multlinegap0pt
\usepackage[all]{xy}
\usepackage{amssymb}
\usepackage{enumerate}
\newtheorem{lemma}{Lemma}
\newcommand\Ext{\operatorname{Ext}}
\newcommand\Hom{\operatorname{Hom}}
\title{On extension modules of finite homological dimension}
\author{Rafael Holanda, Victor H. Jorge-P\'erez, Victor D. Mendoza-Rubio}
\date{Source: arXiv:2504.15224v2 (CC BY 4.0)}
\begin{document}
\maketitle

\noindent\emph{Source and licence.} On extension modules of finite homological dimension. Version arXiv:2504.15224v2 (CC BY 4.0), distributed by arXiv under the Creative Commons Attribution 4.0 International licence, http://creativecommons.org/licenses/by/4.0/, as recorded in the arXiv licence metadata for this exact version and shown on the abstract page https://arxiv.org/abs/2504.15224v2. Full licence notice: \texttt{licenses/arxiv-2504.15224-CC-BY-4.0.txt}.\par\smallskip

\noindent\textit{Editorial scope.} Counted unit: the article's lemma lemswger (source0.tex lines 488--490). The article's complete original proof of it is delivered with it (source0.tex lines 491--506, one \texttt{proof} environment of 16 lines). No mathematics is written, repaired or re-derived: the statement and the proof are the article's own lines, in the article's own notation, and no retained line is substituted or re-flowed.  No further source-local context is needed: every cross-reference of every retained line resolves inside the two delivered spans.  Nothing was omitted from any retained span, so the package carries no omission record.  No retained line contains a citation, so the wrapper carries no bibliography and no bibliography entry.  No retained line contains a figure environment, an image-inclusion command or drawing code, so no image is delivered and the package declares no asset dependency.  Editorial formatting only, in addition: an AMS \texttt{amsart} preamble that restates the article's own declarations and macros used by the retained lines, namely the environments \texttt{lemma} on separate counters, together with the standard packages those lines need.  The retained environments print in this document as Lemma 1 in order of appearance, whereas the article prints the same items with the numbering of its own full document.  The complete native source of the article as distributed in the arXiv e-print for this exact version is bundled unchanged as \texttt{sources/176858/source0.tex}.
\begin{lemma}\label{lemswger}
Let $\mathcal{X}$ be a class of the category of  $R$-modules closed under short exact sequences and isomorphisms. Let $M$ and $N$ be $R$-modules, and $n$ be a non-negative integer. Suppose that $\Ext_R^i(M,N)$ and $N$ belong to $\mathcal{X}$ for all $0\leq i \leq n$. Then $\Hom_R( \Omega^i(M), N) \in \mathcal{X}$ for all $0\leq i \leq n$. 
\end{lemma}

\begin{proof}
    We must show that for all $0\leq j \leq n$, we have  $\Hom_R(\Omega^j(M), N)) \in \mathcal{X}$. We prove this by induction on $j$. For $j=0$, this holds by hypothesis. Now suppose  $0<j\leq n$ and that $\Hom_R(\Omega^{j-1}(M), N))\in\mathcal{X}$. We have an exact sequence
    $\xymatrix@=1em{0 \ar[r] & \Omega^{j}(M) \ar[r] & R^{n_j} \ar[r] & \Omega^{j-1}(M) \ar[r] & 0,}  $
    which yields 
    $$\xymatrix@=1em{
    0\ar[r] & \Hom_R(\Omega^{j-1}(M),N)\ar[r] & N^{n_j}\ar[r] & \Hom_R(\Omega^j(M),N)\ar[r] & \Ext^j_R(M,N)\ar[r] & 0
    }$$
    since $\Ext^1_R(\Omega^{j-1}(M),N)\simeq\Ext^j_R(M,N)$. The result follows by breaking this exact sequence into two short ones and using the fact that $\mathcal{X}$ is closed under short exact sequences and isomorphisms.
    
%    { \small
% \[ 0 \rightarrow \Hom_R(\Omega^{j-1}(M),N) \rightarrow N^{n_j} \rightarrow \Hom_R(\Omega^j(M), N) \rightarrow \Ext_R^1(\Omega^{j-1}(M),N) \cong  
%  \Ext_R^{j}(M,N)\rightarrow 0.\]
% }
 %It induces two exact sequences, 

\end{proof}

\end{document}
